\textbf{(i)} Suppose $fg=1$, where $g=\sum_{j=0}^m b_jx^j$ and $n>0$. Comparing highest coefficients gives $a_nb_m=0$. Inductively, multiply the coefficient equation of degree $n+m-r$ by $a_n^r$. Every term except $a_n^{r+1}b_{m-r}$ vanishes by the preceding identities, so $a_n^{r+1}b_{m-r}=0$ for $0\leq r\leq m$. Since $a_0b_0=1$, this gives $a_n^{m+1}=0$. Now $f-a_nx^n$ is a unit by \exref{1.1}; induction on $n$ proves the claim.

Conversely, a finite sum of nilpotent elements in a commutative ring is nilpotent. Thus $\sum_{i>0}a_ix^i$ is nilpotent, and \exref{1.1} applies.

\textbf{(ii)} If $f^r=0$, then $a_0^r=0$. Also $1+f$ is a unit, so (i) makes the remaining coefficients nilpotent. The converse follows by expanding a sufficiently large power of $f$.

\textbf{(iii)} Choose a nonzero $g=\sum_{j=0}^m b_jx^j$ of least degree with $fg=0$. Then $a_ng$ annihilates $f$ and has degree less than $m$, so $a_ng=0$. Repeating with the remaining coefficients gives $a_ig=0$ for every $i$. Any nonzero coefficient $b_j$ of $g$ therefore satisfies $b_jf=0$. The converse is immediate.

\textbf{(iv)} Every coefficient of $fg$ belongs to the coefficient ideal of each factor. Thus primitivity of $fg$ implies primitivity of both factors. Conversely, if all coefficients of $fg$ lie in a maximal ideal $\mathfrak m$, then the product of the reductions of $f,g$ in $(A/\mathfrak m)[x]$ is zero. This polynomial ring is a domain, so one factor has all coefficients in $\mathfrak m$ and is not primitive.
